\section{نتیجه‌گیری نهایی و راه مقاله}
\label{sec:final}

\subsection{آنچه با اطمینان می‌شود گفت}

\begin{enumerate}
\item \textbf{\lr{EEG} به پیش‌بینی \lr{EMG} چیزی اضافه نمی‌کند.} در سه نسل مدل (\lr{v1}، \lr{v2} و \lr{v3}) هر بار که مبنا کامل‌تر شد، سهم \lr{EEG} کوچک‌تر شد. در \lr{v3}، وقتی حرکت گذشته و زمان رخدادهای گذشته معلوم است، \lr{EEG} در هیچ افقی از صفر تا یک ثانیه و در هیچ فاز تکلیفی خطا را کم نمی‌کند. \lr{EEG} یک \lr{trial} دیگر با همان زمان‌بندی هم همان نتیجه را می‌دهد.
\item \textbf{آنچه \lr{EEG} در این داده می‌داند، بیشتر ساعت تکلیف است.} یک مدل خطی که فقط زمان رخدادها را می‌داند به $r$ برابر \lr{۰٫۸۱۴} می‌رسد و از مدل فقط‌\lr{EEG} ($r$ برابر \lr{۰٫۶۴۵}) در هر ۱۲ نفر بهتر است.
\item \textbf{\lr{EEG} به خطای پیش‌بینی وزن پاسخ می‌دهد، ولی بعد از حرکت.} در پایلوت فصل~\ref{sec:pe-pilot}، پاسخ \lr{EEG} به جسم سبک‌تر از انتظار در سطح گروه دیده شد. ولی همیشه بعد از تفاوت نیرو و \lr{EMG} شروع شد، و انتظار پیش از تماس از \lr{EEG} رمزگشایی نشد.
\end{enumerate}

\subsection{آنچه هنوز باز است}

گام‌های ۷ و ۸ \lr{v3} (ارزیابی \lr{E0} و \lr{Eguided} روی سری‌های آزمون) هنوز نتیجه نداده‌اند. نتیجه‌ی \lr{validation} نشان می‌دهد که احتمالاً تفاوتی نخواهند داشت، ولی این فقط یک پیش‌بینی است. این پرسش ثانویه است و نتیجه‌ی اصلی را تغییر نمی‌دهد.

\subsection{آیا این کار مقاله می‌شود؟}

اگر ادعای مقاله این باشد که \lr{EEG} اطلاعات تازه‌ای درباره‌ی فعالیت عضله دارد، پاسخ منفی است. هیچ‌کدام از تحلیل‌ها چنین ادعایی را پشتیبانی نمی‌کند. نوشتن نتیجه به شکل «مدل ما نتوانست» هم کار ضعیفی است، چون داور می‌تواند بگوید مدل ضعیف بوده.

راه قوی‌تر یک مقاله‌ی \textbf{ممیزی} است. چند کار منتشرشده روی همین داده دقت بالایی از \lr{EEG} گزارش کرده‌اند:
\begin{itemize}
\item رمزگشایی پوش \lr{EMG} با \lr{CNN-LSTM} (\lr{Amiri} و \lr{Shalchyan}، ۲۰۲۵)؛
\item رمزگشایی نیروی گرفتن با $R^2$ حدود \lr{۰٫۸} (\lr{Dangi} و \lr{Meena}، ۲۰۲۵ و ۲۰۲۶).
\end{itemize}
در خلاصه‌ی این کارها مبنای زمان‌بندی تکلیف یا کنترل جابه‌جایی \lr{trial} دیده نمی‌شود. مقاله‌ی ممیزی این روش‌ها را با پروتکل خودشان بازسازی می‌کند. بعد نشان می‌دهد آیا مدلی که فقط زمان تکلیف را می‌داند و \lr{EEG} جابه‌جاشده به همان دقت می‌رسند یا نه. ابزارهای این کار در \lr{v3} ساخته شده‌اند: مدل‌های زمان‌بندی علّی، \lr{phase-matched swap}، کنترل کانال‌های پیرامونی و زمان شروع فعالیت.

این مسیر دو ریسک دارد:
\begin{itemize}
\item باید روش‌های دیگران دقیق بازسازی شوند و متن کامل آن مقاله‌ها هنوز خوانده نشده است.
\item اگر آن روش‌ها با کنترل‌ها هم دوام بیاورند، ادعای مقاله برعکس می‌شود. البته این هم یک نتیجه‌ی گزارش‌کردنی است.
\end{itemize}
مجله‌های مناسب این کار \lr{Journal of Neural Engineering} یا \lr{Imaging Neuroscience} است. نتیجه‌ی پایلوت خطای پیش‌بینی می‌تواند به‌عنوان یک بخش تکمیلی در همین مقاله بیاید: نمونه‌ای دیگر از اینکه پاسخ \lr{EEG} در این تکلیف از حرکت جدا نمی‌شود.

\subsection{پیشنهاد برای تصمیم تیم}

\begin{enumerate}
\item پذیرفتن نتیجه‌ی منفی خط \lr{EEG} به \lr{EMG} و کنار گذاشتن مدل‌های تازه برای این هدف.
\item خواندن متن کامل کارهای منتشرشده روی \lr{WAY-EEG-GAL} و تصمیم درباره‌ی مقاله‌ی ممیزی بر اساس اینکه کدام‌ها قابل بازسازی‌اند.
\item تمام کردن گام‌های ۷ و ۸ \lr{v3} تا پرسش ثانویه هم بسته شود.
\item نگه داشتن پایلوت خطای پیش‌بینی به‌عنوان بخش تکمیلی، نه به‌عنوان مقاله‌ی جدا.
\end{enumerate}
